% ===== overview =====
\begin{frame}{이번 주차 개요}
  \begin{itemize}
    \item \kb{13.1 Random Walk 기초} — 그래프 위 \kb{무작위 걷기}가 마르코프
      체인임을 정의하고, 카라테 클럽에서 걷기 말뭉치를 만들고,
      노드에 머무는 장기 비율이 도(degree)에 비례함을 유도한다.
    \item \kb{13.2 Node2Vec} — 걷기 경로 = ``문장'', 노드 = ``단어''로
      word2vec의 \kb{skip-gram}을 그대로 학습. $p$·$q$ 두 계수가
      탐색 스타일(BFS형 vs DFS형)을 어떻게 바꾸는지 작은 실험으로 본다.
    \item \kb{13.3 PageRank} — ``영원히 돌아다닌 뒤 각 노드에 머무는
      확률은?''이라는 다른 질문. 댐핑 팩터 + 순간이동, 거듭제곱법
      손계산, 그리고 두 가지 흔한 실수.
    \item \kb{13.4 그래프 학습의 확장} — 두 아이디어를 실제 문제로:
      단백질(AlphaFold $\to$ DeepFRI), 금융(AML World), 교통(STGCN).
  \end{itemize}
  \vskip0.5em
  \kq{그래프로 표현된 데이터를, 신경망이나 통계 모델이 다룰 수 있는
      벡터로 어떻게 바꾸는가?}
\end{frame}
